\documentclass[10pt,a4paper]{amsart}
\usepackage[margin=19mm]{geometry}
\usepackage{amsmath,amssymb,mathtools}
\usepackage[scr=rsfs,cal=euler]{mathalfa}
\usepackage{xfrac}
\usepackage[unicode,hidelinks,bookmarks=false]{hyperref}
\newcommand{\ie}{i.e.\ }
\newcommand{\cf}{cf.\ }
\newcommand{\resp}{respectively\ }
\newcommand{\Subsection}{\subsection}
\providecommand{\nobreakdash}{\nobreak\hbox{-}}
\setlength{\emergencystretch}{2em}
\multlinegap0pt
\usepackage{bm}
\usepackage{bbm}
\usepackage{dsfont}
\usepackage{xspace}
\usepackage{cancel}
\usepackage{mathtools}
\usepackage{amsthm}
\makeatletter
\@ifundefined{theorem}{\newtheorem{theorem}{Theorem}}{}
\@ifundefined{lemma}{\newtheorem{lemma}[theorem]{Lemma}}{}
\makeatother
\title[CAT(0) triangle-pentagon complexes]{CAT(0) triangle-pentagon complexes}
\author{Ioana-Claudia Laz\u{a}r}
\date{Source: arXiv:2601.09297v1 (CC BY 4.0)}
\begin{document}
\maketitle

\noindent\emph{Source and licence.} CAT(0) triangle-pentagon complexes. Version arXiv:2601.09297v1 (CC BY 4.0), distributed under the Creative Commons Attribution 4.0 International licence, as recorded in the arXiv licence metadata for this exact version and shown on the abstract page https://arxiv.org/abs/2601.09297v1. Full rights notice: licenses/arxiv-2601.09297-CC-BY-4.0.txt.\par\smallskip

\noindent\textit{Editorial scope.} Counted unit: the Theorem whose source label is \texttt{3.3} in the article (source0.tex lines 328--330, source label \texttt{3.3}), delivered together with the article's own complete original proof of it (source0.tex lines 332--370, one \texttt{proof} environment of 39 lines). No mathematics is written, repaired or re-derived: statement and proof are the article's own text line for line. Required context retained uncounted: 3.1 (lines 272--274); 3.2 (lines 288--290). Every definition, display and label of those lines is retained unchanged. Omitted from this package and not redistributed: 1--271, 275--287, 291--327, 331--331, 371--716. Nothing inside a retained excerpt is omitted or altered: every retained body is delivered exactly as the article's lines are written, so no omission receipt of any kind accompanies this package. Citations: the retained excerpts carry no citation command at all, so no bibliography entry of the article is transcribed and no reference is redistributed. Reference adaptation: none was needed and none was made; every reference inside the delivered unit resolves inside it, so no unresolved reference or citation is produced. Printed slips kept literal: no printed word, symbol, hypothesis or number of the counted statement or of its proof was altered. Layout and class: the article's own class is \texttt{amsart} with packages including amsmath, amsfonts, amssymb, amsthm, url, graphicx, color, enumerate; the wrapper is the lane's pinned \texttt{amsart} editorial wrapper at 10pt on A4, which restates the theorem-like environments the retained text needs and adds the article's own macro definitions that the retained text needs (none), with extra wrapper packages (bm, bbm, dsfont, xspace, cancel, mathtools). The delivered numbering is the unit's own. Assets: no retained span contains a figure environment, a graphics-inclusion command, a PDF-page inclusion, a table or a TikZ picture; no image file is copied into this package, the package declares no asset dependency and no image file is redistributed with it. Licence: version arXiv:2601.09297v1 of this article is distributed under the Creative Commons Attribution 4.0 International licence, as recorded in the arXiv licence metadata for this exact version and shown on the abstract page https://arxiv.org/abs/2601.09297v1. A full rights notice is delivered at licenses/arxiv-2601.09297-CC-BY-4.0.txt. Characters: every retained excerpt is pure ASCII, so the delivered engine, XeLaTeX with its default Latin Modern text font, reports no missing character.
\begin{lemma}\label{3.1}
Let $X$ be a triangle-pentagon complex endowed with a CAT(0) metric $d$. Let $X^{\star}$ be endowed with the metric $d^{\star}$. Then $d^{\star}$ is a CAT(0) metric.
\end{lemma}

\begin{lemma}\label{3.2}
Let $X$ be a CAT(0) triangle-pentagon complex.  Let $X^{\star}$ be endowed with the metric $d^{\star}$.  Then there are no $4$-large and no $3$-large vertices in $X^{\star}$. In particular, $X^{\star}$ is flag and locally $5$-large. Moreover, any $5$-large vertex of $X^{\star}$ is the center of a pentagon of $X$.
\end{lemma}

\begin{theorem}\label{3.3}
Let $X$ be a CAT(0) triangle-pentagon complex.  Let $X^{\star}$ be endowed with the metric $d^{\star}$. Then $X^{\star}$ is $7$-located. 
\end{theorem}

\begin{proof}
Lemma \ref{3.1} implies that $d^{\star}$ is a CAT(0) metric.

We show that any loop $\gamma$ in $X^{\star}$ of length at most $7$ at the boundary of a dwheel, is in the link of a vertex, i.e., there is a vertex $u$ of $X^{\star}$ such that $\gamma \subset X^{\star}_{u}$. 
Due to the Lemma \ref{3.2}, in $X^{\star}$ there are no $3$-large and no $4$-large vertices. 

 
We show that the only $5$-large vertices of $X^{\star}$ are the centers of the pentagons of $X$. Suppose in $X^{\star}$ there exists a $5$-large vertex $v$ which is not the center of a pentagon of $X$. Then there are two cases treated below.

$\bullet$ Suppose $v$ lies at the intersection of a pentagon of $X$ and three equilateral triangles of $X$ (and therefore also of $X^{\star}$). Hence in $X^{\star}$ the sum of the measures of the angles at $v$ is
\begin{center}
$2 \cdot \cfrac{3 \pi}{10} + 3 \cdot \cfrac{\pi}{3} < 2 \pi.$
\end{center} This yields a contradiction with $X^{\star}$ being a CAT(0) space.


$\bullet$ Suppose $v$ lies at the intersection of five equilateral triangles of $X$ (and therefore also of $X^{\star}$). Hence in $X^{\star}$ the sum of the measures of the angles at $v$ is
\begin{center}
$5 \cdot \cfrac{\pi}{3} < 2 \pi.$
\end{center} This yields a contradiction with $X^{\star}$ being a CAT(0) space. 

In conclusion the only $5$-large vertices $v$ of $X^{\star}$ are the centers of pentagons. Let $v$ be the center of the pentagon $\tau = \langle v_{1}, v_{2}, v_{3}, v_{4}, v_{5} \rangle$ of $X$. Then $X_{v} = (v_{1}, v_{2}, v_{3}, v_{4}, v_{5})$.
Let $w=v_{2}$ be a vertex of $X$ adjacent to $v$. Due to Lemma \ref{3.2}, $w$ is not $3$-large and it is not $4$-large.


$\bullet$ Suppose $w$ is $5$-large. Let $X_{w} = (w_{1}=v_{1}, w_{2}, w_{3}, w_{4}=v_{3}, w_{5} = v)$. Then $\gamma = (w_{1}=v_{1}, w_{2}, w_{3}, w_{4}=v_{3}, v_{4}, v_{5})$ is a loop in $X^{\star}$ at the boundary of a dwheel. Note that $|\gamma| = 6$. Note that the sum of the measures of the angles around $w$ is
\begin{center}
$2 \cdot \cfrac{3 \pi}{10} + 3 \cdot \cfrac{\pi}{3} = \cfrac{8 \pi}{5} < 2 \pi$. \end{center} This yields a contradiction with the fact that $X^{\star}$ is a CAT(0) space.


$\bullet$ Suppose $w$ is $6$-large. Let $X_{w} = (w_{1}=v_{1}, w_{2}, w_{3}, w_{4}, w_{5}=v_{3}, w_{6} = v)$. Then $\gamma =(w_{1}=v_{1}, w_{2}, w_{3}, w_{4}, w_{5} = v_{3}, v_{4}, v_{5})$ is a loop in $X^{\star}$ at the boundary of a dwheel. Note that $|\gamma| = 7$. Hence the sum of the measures of the angles around $w$ is
\begin{center}
$2 \cdot \cfrac{3 \pi}{10} + 4 \cdot \cfrac{\pi}{3} = \cfrac{29 \pi}{15} < 2 \pi$.\end{center} This yields a contradiction with the fact that $X^{\star}$ is a CAT(0) space.

In conclusion in $X^{\star}$ there are no loops $\gamma$ of length at most $7$ at the boundary of a dwheel. Any loop in $X^{\star}$ at the boundary of a dwheel and of length at most $7$, is therefore filled with a single vertex. Due to Lemma \ref{3.2}, $X^{\star}$ is flag. In conclusion $X^{\star}$ is $7$-located.

%!! Still need to argue why it can not be filled with 3 %vertices. Analogue 5/9-condition. (if we use the other %definition; not the one with dwheels)


\end{proof}

\end{document}
